\begin{name}
	{\tenchude}
	{\tendethi}
	{\tentruong}
	{\thoigian}
\end{name}
\Opensolutionfile{ans}[ans/ans-De02HT]
%\setcounter{ex}{0}
%%%=========Cau_1=========%%%
% \begin{ex}%[1D2Y2-1]
% 	Tập hợp $M$ có $12$ phần tử. Số tập con gồm $2$ phần tử của $M$ là
% 	\choice
% 	{$12^2$}
% 	{\True $\mathrm{C}_{12}^2$}
% 	{$\mathrm{A}_{12}^{10}$}
% 	{$\mathrm{A}_{12}^2$}
% 	\loigiai{
% 		Số tập con thỏa mãn đề bài chính là số cách chọn $2$ phần tử lấy trong tập hợp $M$ có $12$ phần tử. Số tập con gồm $2$ phần tử của tập hợp $M$ là $\mathrm{C}_{12}^2$.}
% \end{ex}
% %%%=========HetCau_1=========%%%

% %%%=========Cau_2=========%%%
% \begin{ex}%[1D3Y3-3]
% 	Cho cấp số cộng $(u_n)$ có $u_4=-12$ và $ u_{14}=18 $. Giá trị công sai của cấp số cộng đó là
% 	\choice
% 	{$ d=4 $}
% 	{$ d=-3 $}
% 	{\True $ d=3 $}
% 	{$ d=-2 $}
% 	\loigiai{
% 		Ta có $ u_{14}=u_1+13d=u_4+10d=18\Rightarrow d=3 $.\\
% 		Vậy công sai của cấp số cộng là $ d=3 $.}
% \end{ex}
% %%%=========HetCau_2=========%%%

% %%%=========Cau_3=========%%%
% \begin{ex}%[2D1B2-7]
% 	Cho hàm số $ f(x) $ có đạo hàm $ f'(x)=x(x-1)^2(x-2)^5(x-3)^7 $. Số điểm cực trị của hàm số đã cho là
% 	\choice
% 	{\True $ 3 $}
% 	{$ 1 $}
% 	{$ 4 $}
% 	{$ 2 $}
% 	\loigiai{
% 		Ta có $ f'(x)=0\Leftrightarrow x(x-1)^2(x-2)^5(x-3)^7=0\Leftrightarrow \left[\begin{aligned}
% 	& x=0 \\
% 	& x=1 \\
% 	& x=2 \\
% 	& x=3.
% \end{aligned}\right. $ \\
% 		Bảng xét dấu $ f'(x) $ như sau
% 		\begin{center}
% 			\begin{tikzpicture}
% 				\tkzTabInit[nocadre,lgt=1.2,espcl=2.5,deltacl=0.6]
% 				{$ x $ /1, $ f'(x) $ /1}
% 				{$ -\infty $, $ 0 $, $ 1 $, $ 2 $, $ 3 $, $ +\infty $}
% 				\tkzTabLine{,-,0,+,0,+,0,-,0,+,}
% 			\end{tikzpicture}
% 		\end{center}
% 		Từ bảng xét dấu ta thấy $ f'(x) $ có 3 lần đổi dấu nên hàm số đã cho có 3 điểm cực trị.}
% \end{ex}
% %%%=========HetCau_3=========%%%

% %%%=========Cau_4=========%%%
% \begin{ex}%[2D1Y2-2]
% Cho hàm số $ f(x) $ có bảng biến thiên như hình vẽ. 
% \begin{center}
% 	\begin{tikzpicture}[scale=0.7, font=\footnotesize, line join=round, line cap=round, >=stealth]
% 	\tkzTabInit[nocadre,lgt=1.3,espcl=2.5,deltacl=0.6]
% 	{$x$ /1,$f'(x)$ /1, $f(x)$/2 }
% 	{$-\infty$,$1$,$3$,$+\infty$} 
% 	\tkzTabLine{,+,0,-,0,+,} 
% 	\tkzTabVar{-/ $-\infty$,+/$-1 $,-/$-3$,+/$+\infty$}  
% 	\end{tikzpicture}
% 	\end{center}
% 	Điểm cực đại của hàm số đã cho là
% 	\choice 
% 	{$ x=-3 $}
% 	{$ x=3 $}
% 	{$ x=-1 $}
% 	{\True $ x=1 $}
% 	\loigiai{
% 		Hàm số đạt cực đại tại điểm $ x $ mà $ f'(x) $ đổi dấu từ dương sang âm.\\
% 		Từ bảng biến thiên ta thấy hàm số đạt cực đại tại $ x=1 $.}
% \end{ex}
% %%%=========HetCau_4=========%%%

% %%%=========Cau_5=========%%%
% \begin{ex}%[2D1Y4-1]
% 	Tiệm cận ngang của đồ thị hàm số $ y=\dfrac{2x+1}{x-1} $ là
% 	\choice
% 	{$ y=-1 $}
% 	{$ y=1 $}
% 	{$ y=\dfrac{1}{2} $}
% 	{\True $ y=2 $}
% 	\loigiai{
% 		Ta có $ \lim\limits_{x\to \pm \infty}\dfrac{2x+1}{x-1}=\lim\limits_{x\to \pm \infty}\dfrac{2+\dfrac{1}{x}}{1-\dfrac{1}{x}}=2 $.\\ Suy ra đồ thị hàm số có đường tiệm cận ngang là $ y=2 $.}
% \end{ex}
% %%%=========HetCau_5=========%%%

% %%%=========Cau_6=========%%%
% \begin{ex}%[2D1K5-1]
% 	\immini{Đồ thị của hàm số nào dưới đây có dạng như đường cong trong hình vẽ?
% 		\choice
% 		{$ y=-x^4+2x^2 $}
% 		{$ y=x^2-2x+1 $}
% 		{$ y=x^3-3x+1 $}
% 		{\True $ y=-x^3+3x+1 $}}{
% 		\begin{tikzpicture}[scale=0.75, font=\scriptsize, line join=round, line cap=round, >=stealth]
% 			\def\xmin{-2.5}\def\xmax{3}\def\ymin{-2}\def\ymax{4}
% 			\draw[->] (\xmin-0.2,0)--(\xmax+0.2,0) node[below] {$ x $};
% 			\draw[->] (0,\ymin-0.2)--(0,\ymax+0.2) node[right] {$ y $};
% 			\draw (0,0) node [below left] at (0.15,0.1) {$ O $} 
% 			;
% 			\foreach \x in {1}\draw (\x,0.1)--(\x,-0.1) node [below] {$ \x $};
% 			\foreach \x in {-1}\draw (\x,0.1)--(\x,-0.1) node [above] {$ \x $};
% 			\foreach \y in {1,3}\draw (0.1,\y)--(-0.1,\y) node [left] {$ \y $};
% 			\foreach \y in {-1}\draw (0.1,\y)--(-0.1,\y) node [right] {$ \y $};
% 			\clip (\xmin,\ymin) rectangle (\xmax,\ymax);
% 			\draw[smooth,samples=200,domain=\xmin:\xmax] plot (\x,{-1*((\x)^3)+0*((\x)^2)+3*(\x)+1});
% 			\draw[dashed] (1,0)--(1,3)--(0,3) (-1,0)--(-1,-1)--(0,-1);
% 			\fill (1,3) circle (1pt); \fill (-1,-1) circle (1pt); \fill (0,-1) circle (1pt);
% 		\end{tikzpicture}
% 	}
% 	\loigiai{
% 		Đường cong có dạng của đồ thị hàm số bậc $ 3 $ với hệ số $ a<0 $ nên chỉ có hàm số $ y=-x^3+3x+1 $ thỏa yêu cầu bài toán.}
% \end{ex}
% %%%=========HetCau_6=========%%%

% %%%=========Cau_7=========%%%
% \begin{ex}%[2D1B5-4]
% 	\immini{Cho hàm số bậc bốn $ y=f(x) $ có đồ thị là đường cong trong hình vẽ. Số nghiệm của phương trình $ f(x)=-\dfrac{1}{2} $ là
% 		\choice
% 		{\True $ 2 $}
% 		{$ 3 $}
% 		{$ 4 $}
% 		{$ 1 $}}{
% 		\begin{tikzpicture}[scale=0.75, font=\scriptsize, line join=round, line cap=round, >=stealth]
% 			\def\xmin{-2.2}\def\xmax{2.5}\def\ymin{-2.2}\def\ymax{2}
% 			\draw[->] (\xmin-0.2,0)--(\xmax+0.2,0) node[below] {$ x $};
% 			\draw[->] (0,\ymin-0.2)--(0,\ymax+0.2) node[right] {$ y $};
% 			\draw [fill=black] (0,0) node [below left]  at (0.15,0.1) {$ O $}circle(1pt);
% 			\foreach \x in {-2,-1}\draw (\x,0.1)--(\x,-0.1) node [above] {$ \x $};
% 			\foreach \x in {1,2}\draw (\x,0.06)--(\x,-0.06) node [above] {$ \x $};
% 			\foreach \y in {-1,1}\draw (0.1,\y)--(-0.1,\y) node [left] {$ \y $};
% 			\foreach \y in {-2}\draw (0.1,\y)--(-0.1,\y) node [below left] {$ \y $};
% 			\clip (\xmin,\ymin) rectangle (\xmax,\ymax);
% 			\draw[smooth,samples=200,domain=\xmin:\xmax] plot (\x,{1*((\x)^4)+-2*((\x)^2)+-1});
% 			\draw[dashed] (1,0)--(1,-2)--(-1,-2)--(-1,0);
% 			\fill (1,-2) circle (1pt); \fill (-1,-2) circle (1pt); \fill (0,-1) circle (1pt);
% 		\end{tikzpicture}
% 	}
% 	\loigiai{
% 		\immini{Số nghiệm của phương trình $ f(x)=-\dfrac{1}{2} $ bằng số giao điểm của đồ thị hàm số $ y=f(x) $ và đường thẳng $ y=-\dfrac{1}{2}\cdot $\\
% 			Dựa vào đồ thị của hàm số $ y=f(x) $ ta có, đồ thị hàm số $ y=f(x) $ và đường thẳng $ y=-\dfrac{1}{2} $ cắt nhau tại $ 2 $ điểm.\\
% 			Do đó, phương trình $ f(x)=-\dfrac{1}{2} $ có $ 2 $ nghiệm.}{
% 			\begin{tikzpicture}[scale=0.75, font=\scriptsize, line join=round, line cap=round, >=stealth]
% 				\def\xmin{-2.2}\def\xmax{2.5}\def\ymin{-2.2}\def\ymax{2}
% 				\draw[->] (\xmin-0.2,0)--(\xmax+0.2,0) node[below] {$ x $};
% 				\draw[->] (0,\ymin-0.2)--(0,\ymax+0.2) node[right] {$ y $};
% 				\draw [fill=black] (0,0) node [below left] at (0.15,0.1) {$ O $}circle(1pt);
% 				\foreach \x in {-2,-1,1,2}\draw (\x,0.1)--(\x,-0.1) node [above] {$ \x $};
% 				\foreach \y in {-1,1}\draw (0.1,\y)--(-0.1,\y) node [left] {$ \y $};
% 				\foreach \y in {-2}\draw (0.1,\y)--(-0.1,\y) node [below left] {$ \y $};
% 				\draw (\xmin-0.2,-1/2)--(\xmax+0.2,-1/2)node[below]{$ y=-\frac{1}{2} $};
% 				\clip (\xmin,\ymin) rectangle (\xmax,\ymax);
% 				\draw[smooth,samples=200,domain=\xmin:\xmax] plot (\x,{1*((\x)^4)+-2*((\x)^2)+-1});
% 				\draw[dashed] (1,0)--(1,-2)--(-1,-2)--(-1,0);
% 				\fill (1,-2) circle (1pt); \fill (-1,-2) circle (1pt); \fill (0,-1) circle (1pt);
% 			\end{tikzpicture}
% 		}

% 	}
% \end{ex}
% %%%=========HetCau_7=========%%%

% %%%=========Cau_8=========%%%
% \begin{ex}%[2D4Y2-2]
% 	Cho hai số phức $ z_1=5i $ và $ z_2=2020+i $. Phần thực của số $ z_1\cdot z_2 $ bằng
% 	\choice
% 	{\True $ -5 $}
% 	{$ 5 $}
% 	{$ -10~100 $}
% 	{$ 10~100 $}
% 	\loigiai{
% 		Ta có $ z_1\cdot z_2=5i(2020+i)=-5+10~100i $.\\ Phần thực của số phức $ z_1\cdot z_2 $ là $ -5 $.}
% \end{ex}
% %%%=========HetCau_8=========%%%

% %%%=========Cau_9=========%%%
% \begin{ex}%[2D3Y2-1]
% 	Tích phân $ \displaystyle\int\limits_0^1\mathrm{e}^{3x+1}\mathrm{\,d}x $ bằng
% 	\choice
% 	{$ \mathrm{e}^3-\mathrm{e}$}
% 	{$ \dfrac{1}{3}(\mathrm{e}^4+\mathrm{e}) $}
% 	{$ \mathrm{e}^4-\mathrm{e} $}
% 	{\True $ \dfrac{1}{3}(\mathrm{e}^4-\mathrm{e}) $}
% 	\loigiai{
% 		Ta có $ \displaystyle\int\limits_0^1\mathrm{e}^{3x+1}\mathrm{\,d}x=\dfrac{1}{3}\displaystyle\int\limits_0^1\mathrm{e}^{3x+1}\mathrm{\,d}(3x+1)=\dfrac{1}{3}\mathrm{e}^{3x+1}\,\bigg\vert_0^1=\dfrac{1}{3}(\mathrm{e}^4-\mathrm{e}) $.}
% \end{ex}
% %%%=========HetCau_9=========%%%

% %%%=========Cau_10=========%%%
% \begin{ex}%[2H3Y2-4]
% 	Trong không gian với hệ tọa độ $ Oxyz $, cho mặt phẳng $ (P)\colon x-2y+z-5=0 $. Điểm nào dưới đây thuộc $ (P) $? 
% 	\choice
% 	{\True $ M(1;1;6) $}
% 	{$ N(-5;0;0) $}
% 	{$ P(0;0;-5) $}
% 	{$ Q(2;-1;5) $}
% 	\loigiai{
% 		Ta có $ 1-2\cdot1+6-5=0 $ nên $ M(1;1;6) $ thuộc mặt phẳng $ (P) $.}
% \end{ex}
% %%%=========HetCau_10=========%%%

% %%%=========Cau_11=========%%%
% \begin{ex}%[2D2Y4-2]
% 	\renewcommand{\dotEX}{}
% 	Tìm đạo hàm của hàm số $ y=\log_7x $ (với $ x>0$). 
% 	\choice
% 	{$ y'=\dfrac{7}{x}\cdot $}
% 	{$ y'=\dfrac{1}{x}\cdot $}
% 	{\True $ y'=\dfrac{1}{x\ln 7}\cdot $}
% 	{$ y'=\dfrac{\ln 7}{x}\cdot $}
% 	\loigiai{
% 		Đạo hàm của hàm số $ y=\log_7x $ là $ y'=\dfrac{1}{x\ln 7}\cdot $}
% \end{ex}
% %%%=========HetCau_11=========%%%

% %%%=========Cau_12=========%%%
% \begin{ex}%[2H1Y3-2]
% 	\renewcommand{\dotEX}{.}
% 	Cho khối chóp có diện tích đáy $ B=6a^2 $ và chiều cao $ h=2a $. Thể tích khối chóp đã cho bằng
% 	\choice
% 	{$ 12a^3 $}
% 	{$ 2a^3 $}
% 	{\True $ 4a^3 $}
% 	{$ 6a^3 $}
% 	\loigiai{
% 		Ta có $ V=\dfrac{1}{3}B\cdot h=\dfrac{1}{3}\cdot 6a^2\cdot 2a=4a^3 $.}
% \end{ex}
% %%%=========HetCau_12=========%%%

% %%%=========Cau_13=========%%%
% \begin{ex}%[2D3B2-1]
% 	Trong các khẳng định sau, khẳng định nào \textbf{sai}?
% 	\choice
% 	{$ \displaystyle\int\dfrac{1}{x}\mathrm{\,d}x=\ln | x |+C $}
% 	{$ \displaystyle\int x^\mathrm{e}\mathrm{\,d}x=\dfrac{x^{\mathrm{e}+1}}{\mathrm{e}+1}+C $}
% 	{\True $ \displaystyle\int \mathrm{e}^x\mathrm{\,d}x=\dfrac{\mathrm{e}^{x+1}}{x+1}+C $}
% 	{$ \displaystyle\int{\cos 2x\mathrm{\,d}x=\dfrac{1}{2}\sin 2x+C} $}
% 	\loigiai{
% 		Ta có $ \displaystyle\int_{}^{}\mathrm{e}^x\mathrm{\,d}x=\dfrac{\mathrm{e}^{x+1}}{x+1}+C $ sai vì $ \displaystyle\int_{}^{}\mathrm{e}^x\mathrm{\,d}x=\mathrm{e}^x+C $.}
% \end{ex}
% %%%=========HetCau_13=========%%%

% %%%=========Cau_14=========%%%
% \begin{ex}
% 	Trong không gian với hệ trục $ Oxyz $, cho $ \vec{a}=(-2;2;0) $, $ \vec{b}=(2;2;0) $, $ \vec{c}=(2;2;2) $. Giá trị của $ \left|\vec{a}+\vec{b}+\vec{c}\, \right| $ bằng
% 	\choice
% 	{$ 2\sqrt{6} $}
% 	{$ 11 $}
% 	{\True $ 2\sqrt{11} $}
% 	{$ 6 $}
% 	\loigiai{
% 		Ta có $ \vec{a}+\vec{b}+\vec{c}=(2;6;2) $,
% 		suy ra $ \left| \vec{a}+\vec{b}+\vec{c}\, \right|=2\sqrt{11} $.}
% \end{ex}
% %%%=========HetCau_14=========%%%

% %%%=========Cau_15=========%%%
% \begin{ex}%[2D2Y5-1]
% 	Phương trình $ 3^{x^2-2x}=1 $ có nghiệm là
% 	\choice
% 	{\True $ x=0 $, $ x=2 $}
% 	{$ x=-1 $, $ x=3 $}
% 	{$ x=0 $, $ x=-2 $}
% 	{$ x=1 $, $ x=-3 $}
% 	\loigiai{
% 		Ta có $ 3^{x^2-2x}=1\Leftrightarrow 3^{x^2-2x}=3^0\Leftrightarrow x^2-2x=0\Leftrightarrow \left[\begin{aligned}
% 	& x=0 \\
% 	& x=2.
% \end{aligned}\right. $}
% \end{ex}
% %%%=========HetCau_15=========%%%

% %%%=========Cau_16=========%%%
% \begin{ex}%[2H3Y3-1]
% 	Trong không gian với hệ trục tọa độ $ Oxyz, $ cho đường thẳng $ d\colon\dfrac{x-3}{1}=\dfrac{y+1}{-2}=\dfrac{z-5}{3} $. Véc-tơ nào sau đây là một véc-tơ chỉ phương của đường thẳng $ d $? 
% 	\choice
% 	{\True $ \vec{u}_2=(1;-2;3) $}
% 	{$ \vec{u}_4=(-2;-4;6) $}
% 	{$ \vec{u}_3=(2;6;-4) $}
% 	{$ \vec{u}_1=(3;-1;5) $}
% 	\loigiai{
% 		Ta thấy đường thẳng $ d $ có một véc-tơ chỉ phương có tọa độ $ \vec{u}_2=(1;-2;3) $.}
% \end{ex}
% %%%=========HetCau_16=========%%%

% %%%=========Cau_17=========%%%
% \begin{ex}%[2D4Y1-2]
% 	\immini{Trong mặt phẳng $ Oxy $, số phức $ z=-2+4i $ được biểu diễn bởi điểm nào trong các điểm ở hình vẽ?
% 		\choice
% 		{\True Điểm $ C $}
% 		{Điểm $ D $}
% 		{Điểm $ A $}
% 		{Điểm $ B $}}{
% 		\begin{tikzpicture}[>=stealth,line join=round,line cap=round,font=\scriptsize,scale=0.45]
% 			\draw [->] (-5,0)--(5,0)node[below]{$ x $};
% 			\draw [->] (0,-5)--(0,5)node[left]{$ y $};
% 			\draw [dashed] (-4,0)node[below]{$ -4 $}--(-4,2)--(0,2)node[right]{$ 2 $} (-2,0)node[below]{$ -2 $}--(-2,4)--(0,4)node[right]{$ 4 $} (4,0)node[above]{$ 4 $}--(4,-2)--(0,-2)node[left]{$ -2 $} (2,0)node[above]{$ 2 $}--(2,-4)--(0,-4)node[left]{$ -4 $};
% 			\draw [fill=black] (0,0)node[below left]{$ O $}circle(1pt) (-4,2)node[left]{$ A $}circle(1pt) (-2,4)node[left]{$ C $}circle(1pt) (2,-4)node[right]{$ B $}circle(1pt) (4,-2)node[right]{$ D $}circle(1pt);
% 		\end{tikzpicture}
% 	}
% 	\loigiai{
% 		Số phức $ z=-2+4i $ được biểu diễn bởi điểm $ C(-2;4) $.}
% \end{ex}
% %%%=========HetCau_17=========%%%

% %%%=========Cau_18=========%%%
% \begin{ex}%[2D3B2-1]
% 	Cho hàm số $ f(x) $ liên tục trên $ \mathbb{R} $ và thỏa mãn $ \displaystyle\int\limits_0^1f(x)\mathrm{\,d}x=2 $, $ \displaystyle\int\limits_1^3f(x)\mathrm{\,d}x=6 $. Tính $ I=\displaystyle\int\limits_0^3f(x)\mathrm{\,d}x $.
% 	\choice
% 	{\True $ I=8 $}
% 	{$ I=12 $}
% 	{$ I=4 $}
% 	{$ I=36 $}
% 	\loigiai{
% 		Ta có $ I=\displaystyle\int\limits_0^3f(x)\mathrm{\,d}x=\displaystyle\int\limits_0^1f(x)\mathrm{\,d}x+\displaystyle\int\limits_1^3f(x)\mathrm{\,d}x=2+6=8 $.}
% \end{ex}
% %%%=========HetCau_18=========%%%

% %%%=========Cau_19=========%%%
% \begin{ex}%[2H1Y3-2]
% 	Khối nón có chiều cao bằng $ 4 $ và đường kính đáy bằng $ 6 $. Thể tích khối nón bằng	
% 	\choice
% 	{\True $ 12\pi $}
% 	{$ 144\pi $}
% 	{$ 48\pi $}
% 	{$ 24\pi $}
% 	\loigiai{
% 		Khối nón có bán kính bằng $ 3 $ và chiều cao bằng $ 4 $ nên có thể tích là \\ $V=\dfrac{1}{3}\pi r^2h=\dfrac{1}{3}\cdot\pi\cdot{3^2}\cdot4=12\pi$.}
% \end{ex}
% %%%=========HetCau_19=========%%%

% %%%=========Cau_20=========%%%
% \begin{ex}%[2H1Y3-2]
% 	Cho khối hộp hình chữ nhật có ba kích thước $ 2 $; $ 4 $; $ 6 $. Thể tích của khối hộp đã cho bằng
% 	\choice
% 	{$ 8 $}
% 	{$ 16 $}
% 	{\True $ 48 $}
% 	{$ 12 $}
% 	\loigiai{
% 		Thể tích của khối hộp đã cho bằng $ 2\cdot 4\cdot 6=48 $.}
% \end{ex}
% %%%=========HetCau_20=========%%%

% %%%=========Cau_21=========%%%
% \begin{ex}%[2D4Y2-1]
% 	Cho hai số phức $ z_1=1-2i $ và $ z_2=2+i $. Số phức $ z_1+z_2 $ bằng
% 	\choice
% 	{$ -3-i $}
% 	{$ 3+i $}
% 	{\True $ 3-i $}
% 	{$ -3+i $}
% 	\loigiai{
% 		Ta có $ z_1+z_2=1-2i+2+i=3-i $.}
% \end{ex}
% %%%=========HetCau_21=========%%%

% %%%=========Cau_22=========%%%
% \begin{ex}%[2H3Y1-3]
% 	Trong không gian với hệ trục $ Oxyz $, cho mặt cầu $ (S)\colon x^2+y^2+z^2-4x+2y-6z+1=0 $. Tọa độ tâm $ I $ của mặt cầu là
% 	\choice
% 	{$ I(4;-2;6) $}
% 	{\True $ I(2;-1;3) $}
% 	{$ I(-4;2;-6) $}
% 	{$ I(-2;1;-3) $}
% 	\loigiai{
% 		Từ phương trình mặt cầu suy ra tâm của mặt cầu là $ I(2;-1;3) $.}
% \end{ex}
% %%%=========HetCau_22=========%%%

% %%%=========Cau_23=========%%%
% \begin{ex}%[2D1Y1-2]
% 	Cho hàm số $ y=f(x) $ có bảng biến thiên như hình vẽ.	
% 	\begin{center}
% 			\begin{tikzpicture}[scale=0.7]
% 			\tkzTabInit[nocadre,lgt=1.2,espcl=2.7,deltacl=0.7]
% 			{$ x $ /1, $ y' $ /1, $ y $ /2}
% 			{$ -\infty $, $ -1 $, $ 0 $, $ 1 $, $ +\infty $}
% 			\tkzTabLine{,+,0,-,d,-,0,+,}
% 			\tkzTabVar{-/ $ -\infty $,+/ $ 2 $,-D+/ $ -\infty $ / $ +\infty $,-/ $ 4 $,+/ $ +\infty $}
% 		\end{tikzpicture}
% 	\end{center}
% 	Hàm số nghịch biến trong khoảng nào?
% 		\choice
% 		{\True $ (0;1) $}
% 		{$ (-1; 1) $}
% 		{$ (4;+\infty) $}
% 		{$ (-\infty;2) $}
% 	\loigiai{
% 		Từ bảng biến thiên ta thấy hàm số đã cho nghịch biến trên khoảng $ (0;1) $.}
% \end{ex}
% %%%=========HetCau_23=========%%%

% %%%=========Cau_24=========%%%
% \begin{ex}%[2D2Y5-1]
% 	Nghiệm của phương trình $ \log_2(x+9)=5 $ là
% 	\choice
% 	{$ x=41 $}
% 	{$ x=16 $}
% 	{\True $ x=23 $}
% 	{$ x=1 $}
% 	\loigiai{
% 		Điều kiện $ x>-9 $.
% 		Ta có $ \log_2(x+9)=5\Leftrightarrow x+9=2^5\Leftrightarrow x=23 $.}
% \end{ex}
% %%%=========HetCau_24=========%%%

% %%%=========Cau_25=========%%%
% \begin{ex}%[2D2Y1-1]
% 	Cho $ x, y>0 $ và $ \alpha, \beta \in \mathbb{R} $. Khẳng định nào sau đây \textbf{sai}?
% 	\choice
% 	{$ (x^{\alpha})^{\beta}=x^{\alpha \beta} $}
% 	{\True $ x^{\alpha}+y^{\alpha}=(x+y)^{\alpha} $}
% 	{$ x^{\alpha}\cdot x^{\beta}=x^{\alpha +\beta} $}
% 	{$ (x y)^{\alpha}=x^{\alpha}\cdot y^{\alpha} $}
% 	\loigiai{
% 		Theo tính chất của lũy thừa thì đẳng thức $ x^{\alpha}+y^{\alpha}=(x+y)^{\alpha} $ sai.}
% \end{ex}
% %%%=========HetCau_25=========%%%

% %%%=========Cau_26=========%%%
% \begin{ex}%[2H1Y3-1]
% 	Cho hình trụ có bán kính đáy $ r=2 $ và chiều cao $ h=5 $. Diện tích xung quanh của hình trụ đã cho bằng
% 	\choice
% 	{$ 28\pi $}
% 	{$ 20 $}
% 	{$ 10\pi $}
% 	{\True $ 20\pi $}
% 	\loigiai{
% 		Theo công thức tính diện tích xung quanh hình trụ, ta có\\ $ S_{xq}=2\pi rh=2\pi\cdot 2\cdot 5=20\pi $.}
% \end{ex}
% %%%=========HetCau_26=========%%%

% %%%=========Cau_27=========%%%
% \begin{ex}%[2H3K3-7]
% 	\renewcommand{\dotEX}{}
% 	Trong không gian với hệ trục tọa độ $ Oxyz $, cho các điểm $ A(1;0;2) $, $ B(1;2;1) $, $ C(3;2;0) $ và $ D(1;1;3) $. Đường thẳng đi qua $ A $ và vuông góc với mặt phẳng $ (BCD) $ có phương trình là
% 	\choice
% 	{$ \left\{\begin{aligned}
% 	& x=1-t \\
% 	& y=4t \\
% 	& z=2+2t.
% \end{aligned}\right. $}
% 	{$ \left\{\begin{aligned}
% 	& x=1+t \\
% 	& y=4 \\
% 	& z=2+2t.
% \end{aligned}\right. $}
% 	{$ \left\{\begin{aligned}
% 	& x=1-t \\
% 	& y=2-4t \\
% 	& z=2-2t.
% \end{aligned}\right. $}
% 	{\True $ \left\{\begin{aligned}
% 	& x=2+t \\
% 	& y=4+4t \\
% 	& z=4+2t.
% \end{aligned}\right. $}
% 	\loigiai{
% 		Đường thẳng đi qua $ A $ và vuông góc với mặt phẳng $ (BCD) $ nhận véc-tơ pháp tuyến của $ (BCD) $ làm véc-tơ chỉ phương.\\
% 		Ta có $ \overrightarrow{BC}=(2;0;-1) $, $ \overrightarrow{BD}=(0;-1;2) $ 
% 		 $ \Rightarrow \vec{u}_d=\vec{n}=\left[\overrightarrow{BC},\overrightarrow{BD}\right]=(-1;-4;-2) $.\\ 
% 		Khi đó, ta loại phương án $ \left\{\begin{aligned}
% 	& x=1-t \\
% 	& y=4t \\
% 	& z=2+2t
% \end{aligned}\right. $ và $ \left\{\begin{aligned}
% 	& x=1+t \\
% 	& y=4 \\
% 	& z=2+2t.
% \end{aligned}\right. $ \\
% 		Thay điểm $ A(1;0;2) $ vào phương trình ở phương án $ \left\{\begin{aligned}
% 	& x=2+t \\
% 	& y=4+4t \\
% 	& z=4+2t
% \end{aligned}\right. $, ta có $ \left\{\begin{aligned}
% 	& 1=2+t \\
% 	& 0=4+4t \\
% 	& 2=4+2t
% \end{aligned}\right.\Leftrightarrow \left\{\begin{aligned}
% 	& t=-1 \\
% 	& t=-1 \\
% 	& t=-1.
% \end{aligned}\right. $ \\
% 		Suy ra đường thẳng có phương trình tham số $ \left\{\begin{aligned}
% 	& x=2+t \\
% 	& y=4+4t \\
% 	& z=4+2t
% \end{aligned}\right. $ đi qua điểm $ A $ nên $ \left\{\begin{aligned}
% 	& x=2+t \\
% 	& y=4+4t \\
% 	& z=4+2t
% \end{aligned}\right. $ là phương án đúng.
% 	}
% \end{ex}
% %%%=========HetCau_27=========%%%
% 	\renewcommand{\dotEX}{.}
% %%%=========Cau_28=========%%%
% \begin{ex}%[2D2Y1-2]
% 	Rút gọn biểu thức $ P=\dfrac{a^{\sqrt{3}+1}\cdot a^{2-\sqrt{3}}}{\left(a^{\sqrt{2}-2}\right)^{\sqrt{2}+2}} $ với $ a>0 $.
% 	\choice
% 	{$ P=a^4 $}
% 	{$ P=a^3 $}
% 	{\True $ P=a^5 $}
% 	{$ P=a $}
% 	\loigiai{
% 		Ta có $ P=\dfrac{a^{\sqrt{3}+1}\cdot a^{2-\sqrt{3}}}{\left(a^{\sqrt{2}-2}\right)^{\sqrt{2}+2}}=\dfrac{a^{\sqrt{3}+1+2-\sqrt{3}}}{a^{\left(\sqrt{2}-2\right)\left(\sqrt{2}+2\right)}}=\dfrac{a^3}{a^{-2}}=a^5 $.}
% \end{ex}
% %%%=========HetCau_28=========%%%

% %%%=========Cau_29=========%%%
% \begin{ex}%[2D3B2-1]
% 	Cho $ \displaystyle\int\limits_0^1f(x)\mathrm{\,d}x=2 $ và $ \displaystyle\int\limits_0^1g(x)\mathrm{\,d}x=5 $. Tính $ \displaystyle\int\limits_0^1[f(x)-2g(x)]\mathrm{\,d}x $.
% 	\choice
% 	{\True $ -8 $}
% 	{$ 12 $}
% 	{$ 1 $}
% 	{$ -3 $}
% 	\loigiai{
% 		Ta có $ \displaystyle\int\limits_0^1[f(x)-2g(x)]\mathrm{\,d}x=\displaystyle\int\limits_0^1f(x)\mathrm{\,d}x-2\displaystyle\int\limits_0^1 g(x)\mathrm{\,d}x=2-2\cdot 5=-8 $.}
% \end{ex}
% %%%=========HetCau_29=========%%%

% %%%=========Cau_30=========%%%
% \begin{ex}%[1H3B5-3]
% 	\renewcommand{\dotEX}{}
% 	Cho hình chóp $ S.ABCD $ có $ SA\perp (ABCD) $, đáy $ ABCD $ là hình chữ nhật. Biết $ AD=2a $, $ SA=a $. Khoảng cách từ $ A $ đến $ (SCD) $ bằng
% 	\choice
% 	{$ \dfrac{3a}{\sqrt{7}}\cdot $}
% 	{$ \dfrac{3a\sqrt{2}}{2}\cdot $}
% 	{\True $ \dfrac{2a}{\sqrt{5}}\cdot $}
% 	{$ \dfrac{2a\sqrt{3}}{3}\cdot $}
% 	\loigiai{
% 		\immini{Gọi $ H $ là hình chiếu của $ A $ lên $ SD $, ta chứng minh được $ AH\perp (SCD) $ hay khoảng cách từ $ A $ đến $ (SCD) $ là $ AH $.\\
% 			Ta có $ \dfrac{1}{AH^2}=\dfrac{1}{SA^2}+\dfrac{1}{AD^2}\Rightarrow AH=\dfrac{2a}{\sqrt{5}}\cdot $\\
% 		Vậy $ \mathrm{d}[A,(SCD)]=\dfrac{2a}{\sqrt{5}}\cdot $}{
% 			\begin{tikzpicture}[>=stealth,line join=round,line cap=round,font=\scriptsize,scale=0.65]
% 				\coordinate (A) at (0,0);
% 				\coordinate (B) at (-2,-1);
% 				\coordinate (D) at (4,0);
% 				\coordinate (C) at ($ (B)+(D)-(A) $);
% 				\coordinate (S) at ($ (A)+(0,4) $);
% 				\coordinate (H) at ($ (S)!0.4!(D) $);
% 				\draw (S)--(B)--(C)--(D)--(S)--(C);
% 				\draw [dashed] (B)--(A)--(S) (H)--(A)--(D);
% 				\draw [fill=black]  (B)node[below left]{$ B $}circle(1pt) (C)node[below right]{$ C $}circle(1pt) (D)node[above right]{$ D $}circle(1pt) (S)node[above left]{$ S $}circle(1pt);
% 				\foreach \x/\dinh/\y in {S/A/B,A/H/S}\draw[fill = gray!50] ($ (\dinh)!5pt!(\x) $)--($ (\dinh)!5pt!(\x)+(\dinh)!5pt!(\y)-(\dinh) $)--($ (\dinh)!5pt!(\y) $)--(\dinh)--cycle; 
% 				\draw [fill=black] (A)node[above left]{$ A $}circle(1pt) (H)node[above right]{$ H $}circle(1pt);
% 			\end{tikzpicture}
% 		}
% 	}
% \end{ex}
% %%%=========HetCau_30=========%%%

% %%%=========Cau_31=========%%%
% \begin{ex}%[2D1B3-1]
% 	Giá trị nhỏ nhất của hàm số $ y=x^3+3x^2 $ trên đoạn $ [-4;-1] $ bằng
% 	\choice
% 	{$ 0 $}
% 	{$ 4 $}
% 	{\True $ -16 $}
% 	{$ -4 $}
% 	\loigiai{
% 		Ta có $ y'=3x^2+6x $;\\
% 		 $ y'=0\Rightarrow 3x^2+6x=0\Leftrightarrow \left[\begin{aligned}
% 	& x=0\notin [-4;-1 ] \\
% 	& x=-2\in [-4;-1 ].
% \end{aligned}\right. $ \\
% 		Khi đó, $ y(-4)=-16 $, $ y(-2)=4 $, $ y(-1)=2 $.\\
% 		Do đó, $ \min\limits_{[-4;-1 ]} y=-16 $.}
% \end{ex}
% %%%=========HetCau_31=========%%%

% %%%=========Cau_32=========%%%
% \begin{ex}%[1D2B5-2]
% 	\renewcommand{\dotEX}{}
% 	Một em bé có bộ $ 6 $ thẻ chữ, trên mỗi thẻ có ghi một chữ cái, trong đó có $ 3 $ thẻ chữ $ \boldsymbol{T} $, một thẻ chữ $ \boldsymbol{N} $, một thẻ chữ $ \boldsymbol{H} $ và một thẻ chữ $ \boldsymbol{P} $. Em bé đó xếp ngẫu nhiên $ 6 $ thẻ đó thành một hàng ngang. Tính xác suất em bé xếp được thành dãy $ \boldsymbol{TNTHPT} $.\\
% 	\choice
% 	{\True $ \dfrac{1}{120}\cdot $}
% 	{$ \dfrac{1}{720}\cdot $}
% 	{$ \dfrac{1}{6}\cdot $}
% 	{$ \dfrac{1}{20}\cdot $}
% 	\loigiai{
% 		Xem ba chữ $ \boldsymbol{T} $ riêng biệt ta có $ n(\Omega)=6! $.\\
% 		Gọi $ A $ là biến cố \lq\lq xếp ngẫu nhiên 6 thẻ đó thành dãy $ \boldsymbol{TNTHPT} $\rq\rq, suy ra $ n(A)=3! $.\\
% 		(số hoán vị của $ \boldsymbol{T-T-T} $ và $ \boldsymbol{N} $, $ \boldsymbol{H} $, $ \boldsymbol{P} $ cố định).\\
% 		Vậy xác suất của biến cố $ A $ là $ P(A)=\dfrac{3!}{6!}=\dfrac{1}{120}\cdot $}
% \end{ex}
% %%%=========HetCau_32=========%%%

% %%%=========Cau_33=========%%%
% \begin{ex}%[2D3B1-1]
% 	Tính $ \displaystyle\int{(x-\sin 2x)\mathrm{\,d}x} $ 
% 	\choice
% 	{$ x^2+\dfrac{\cos 2x}{2}+C $}
% 	{\True $ \dfrac{x^2}{2}+\dfrac{\cos 2x}{2}+C $}
% 	{$ \dfrac{x^2}{2}+\cos 2x+C $}
% 	{$ \dfrac{x^2}{2}+\sin x+C $}
% 	\loigiai{
% 		Ta có $ \displaystyle\int (x-\sin 2x)\mathrm{\,d}x=\displaystyle\int x\mathrm{\,d}x-\displaystyle\int \sin 2x\mathrm{\,d}x=\dfrac{x^2}{2}+\dfrac{\cos 2x}{2}+C $.}
% \end{ex}
% %%%=========HetCau_33=========%%%

% %%%=========Cau_34=========%%%
% \begin{ex}%[2D4B3-2]
% 	Cho số phức $ z $ thỏa mãn điều kiện $ (1+i)\overline{z}-1-3i=0 $. Tìm phần ảo của số phức $ w=1-iz+\overline{z} $ 
% 	\choice
% 	{\True $ -1 $}
% 	{$ -i $}
% 	{$ 2 $}
% 	{$ -2i $}
% 	\loigiai{
% 		Ta có $ (1+i)\overline{z}-1-3i=0\Leftrightarrow \overline{z}=\dfrac{1+3i}{1+i}\Leftrightarrow \overline{z}=2+i\Rightarrow z=2-i $.\\
% 		Do đó $ w=1-iz+\overline{z}=1-i(2-i)+2+i=2-i $.\\
% 		Vậy phần ảo của số phức $ w=1-iz+\overline{z} $ là $ -1 $.}
% \end{ex}
% %%%=========HetCau_34=========%%%

% %%%=========Cau_35=========%%%
% \begin{ex}%[2H3Y1-3]
% 	Trong không gian với hệ trục tọa độ $ Oxyz $, cho hai điểm $ I(1;1;1) $ và $ A(1;2;3) $. Phương trình mặt cầu có tâm $ I $ và đi qua $ A $ là
% 	\choice
% 	{$ (x+1)^2+(y+1)^2+(z+1)^2=29 $}
% 	{$ (x-1)^2+(y-1)^2+(z-1)^2=25 $}
% 	{\True $ (x-1)^2+(y-1)^2+(z-1)^2=5 $}
% 	{$ (x+1)^2+(y+1)^2+(z+1)^2=5 $}
% 	\loigiai{
% 		Ta có $ IA=\sqrt{(1-1)^2+(2-1)^2+(3-1)^2}=\sqrt{5} $.\\
% 		Mặt cầu tâm $ I $ và đi qua điểm $ A $ nên có bán kính là $ R=IA $.\\ Phương trình mặt cầu cần tìm là\\
% 		 $ (x-x_I)^2+(y-y_I)^2+(z-z_I)^2=R^2\Rightarrow {(x-1)^2}+(y-1)^2+(z-1)^2=5 $.}
% \end{ex}
% %%%=========HetCau_35=========%%%

% %%%=========Cau_36=========%%%
% \begin{ex}%[2D2B6-2]
% 	Số nghiệm nguyên của bất phương trình $ \left(\dfrac{1}{3}\right)^{2x^2-3x-7}>3^{2x-21} $ là
% 	\choice
% 	{\True $ 7 $}
% 	{$ 6 $}
% 	{vô số}
% 	{$ 8 $}
% 	\loigiai{
% 		Ta có $ \begin{aligned}[t]
% 			\left(\dfrac{1}{3}\right)^{2x^2-3x-7}>3^{2x-21}&\Leftrightarrow{3^{-(2x^2-3x-7)}}>3^{2x-21}\\&\Leftrightarrow -(2x^2-3x-7)>2x-21\\&\Leftrightarrow -2x^2+3x+7>2x-21\\&\Leftrightarrow -2x^2+x+28>0\\&\Leftrightarrow -\dfrac{7}{2}<x<4.
% 		\end{aligned} $ \\
% 		Do $ x\in \mathbb{Z} $ nên $ x\in \{-3;-2;-1;0;1;2;3 \} $.\\
% 		Vậy bất phương trình đã cho có $ 7 $ nghiệm nguyên.}
% \end{ex}
% %%%=========HetCau_36=========%%%

% %%%=========Cau_37=========%%%
% \begin{ex}%[2D1B1-1]
% 	Hàm số $ y=\dfrac{2}{3x^2+1} $ nghịch biến trên khoảng nào dưới đây?
% 	\choice
% 	{$ (-1;1) $}
% 	{$ (-\infty;0) $}
% 	{$ (-\infty;+\infty) $}
% 	{\True $ (0;+\infty) $}
% 	\loigiai{
% 		Tập xác định $ \mathscr{D}=\mathbb{R} $.\\ Ta có
% 		 $ y'=\dfrac{-12x}{(3x^2+1)^2}\cdot $\\
% 		Mà $ y'<0\Leftrightarrow x>0 $ nên hàm số $ y=\dfrac{2}{3x^2+1} $ nghịch biến trên khoảng $ (0;+\infty) $.}
% \end{ex}
% %%%=========HetCau_37=========%%%

% %%%=========Cau_38=========%%%
% \begin{ex}%[2D1G3-1]
% 	\immini{Cho hàm số $ f(x) $. Biết hàm số $ f'(x) $ có đồ thị như hình bên. Trên $ [-4; 3 ] $, hàm số $ g(x)=2f(x)+(1-x)^2 $ đạt giá trị nhỏ nhất tại điểm nào?
% 		\choice
% 		{\True $ x=-1 $}
% 		{$ x=3 $}
% 		{$ x=-4 $}
% 		{$ x=-3 $}}{
% 		\begin{tikzpicture}[>=stealth,line join=round, line cap=round,scale=0.65,font=\scriptsize]
% 			\draw[thin,color=gray!50] (-5,-2.5) grid (4,6.5);
% 			\draw[->] (-5,0)--(4,0) node [below]{$ x $};
% 			\draw[->] (0,-2.5)--(0,6.5) node [right]{$ y $};
% 			\node at (0,0) [above right]{$ O $};
% 			\draw[dashed] (-4,0)--(-4,5)--(0,5);
% 			\draw[dashed] (-3,0)--(-3,3)--(0,3);
% 			\draw[dashed] (-1,0)--(-1,2)--(0,2);
% 			\draw[dashed] (3,0)--(3,-2)--(0,-2);
			
% 			\fill (-4,0)node[below]{$ -4 $}circle(1pt) (-3,0)node[below]{$ -3 $}circle(1pt)  (-2,0)node[below]{$ -2 $}circle(1pt) (-1,0)node[below]{$ -1 $}circle(1pt) (0,-2)node[left]{$ -2 $}circle(1pt) (0,2)node[right]{$ 2 $}circle(1pt) (0,3)node[right]{$ 3 $}circle(1pt) (3,0)node[above]{$ 3 $}circle(1pt) (0,5)node[right]{$ 5 $}circle(1pt)
% 			(-4,5.5)node[right]{$ y=f'(x) $};
% 			\draw[samples=5000,smooth] (-4.1,5.5)--(-4,5) to[out=-75,in=160](-3,3)
% 			to[out=-20,in=150] (-1,2)
% 			to[out=-30,in=120] (2,1)
% 			to[out=-60,in=105](3,-2)--(3.1,-2.4);
% 			\draw [fill=black] (-4,5)circle(1pt) (-3,3)circle(1pt) (-1,2)circle(1pt) (3,-2)circle(1pt);
% 		\end{tikzpicture}
% 	}
% 	\loigiai{
% 		Xét hàm số $ g(x)=2f(x)+(1-x)^2 $ trên $ [-4;3 ] $.\\
% 		Ta có $ g'(x)=2\cdot f'(x)-2(1-x) $, 
% 		 $ g'(x)=0\Leftrightarrow f'(x)=1-x $.\\ Trên đồ thị hàm số $ f'(x) $ ta vẽ thêm đường thẳng $ y=1-x $ 
% 		\begin{center}
% 			\begin{tikzpicture}[>=stealth,line join=round, line cap=round,scale=0.65,font=\scriptsize]
% 				\draw[thin,color=gray!50] (-5,-2.5) grid (4,6.5);
% 				\draw[->] (-5,0)--(4,0) node [below]{$ x $};
% 				\draw[->] (0,-2.5)--(0,6.5) node [right]{$ y $};
% 				\node at (0,0) [above right]{$ O $};
% 				\draw[dashed] (-4,0)--(-4,5)--(0,5);
% 				\draw[dashed] (-3,0)--(-3,3)--(0,3);
% 				\draw[dashed] (-1,0)--(-1,2)--(0,2);
% 				\draw[dashed] (3,0)--(3,-2)--(0,-2);
				
% 				\fill (-4,0)node[below]{$ -4 $}circle(1pt) (-3,0)node[below]{$ -3 $}circle(1pt)  (-2,0)node[below]{$ -2 $}circle(1pt) (-1,0)node[below]{$ -1 $}circle(1pt) (0,-2)node[left]{$ -2 $}circle(1pt) (0,2)node[right]{$ 2 $}circle(1pt) (0,3)node[right]{$ 3 $}circle(1pt) (3,0)node[above]{$ 3 $}circle(1pt) (0,5)node[right]{$ 5 $}circle(1pt)
% 				(-4,5.5)node[right]{$ y=f'(x) $};
% 				\draw[samples=5000,smooth] (-4.1,5.5)--(-4,5) to[out=-75,in=160](-3,3)
% 				to[out=-20,in=150] (-1,2)
% 				to[out=-30,in=120] (2,1)
% 				to[out=-60,in=105](3,-2)--(3.1,-2.4);
% 				\draw [fill=black] (-4,5)circle(1pt) (-3,3)circle(1pt) (-1,2)circle(1pt) (3,-2)circle(1pt);
% 				\draw [samples=5000,smooth,domain=-4.2:3.2] plot(\x,{1-(\x)});
% 			\end{tikzpicture}
% 		\end{center}
% 		Từ đồ thị ta thấy $ f'(x)=1-x\Leftrightarrow \left[\begin{aligned}
% 	& x=-4 \\
% 	& x=-1 \\
% 	& x=3.
% \end{aligned}\right. $\\
% 		Bảng biến thiên của hàm số $ g(x) $ như sau
% 		\begin{center}
% 			\begin{tikzpicture}
% 				\tkzTabInit[nocadre,lgt=1.2,espcl=3,deltacl=0.7]
% 				{$ x $ /0.6, $ g'(x) $ /0.6, $ g(x) $ /2}
% 				{$ -4 $, $ -1 $, $ 3 $}
% 				\tkzTabLine{0,-,0,+,0}
% 				\tkzTabVar{+/ $ g(-4) $,-/ $ g(-1) $,+/ $ g(3) $}
% 			\end{tikzpicture}
% 		\end{center}
% 		Vậy $ \min\limits_{[-4;3 ]} g(x)=g(-1)\Leftrightarrow x=-1 $.}
% \end{ex}
% %%%=========HetCau_38=========%%%

% %%%=========Cau_39=========%%%
% \begin{ex}%[2D1T3-6]
% 	Người ta muốn xây bể chứa nước dạng hình hộp chữ nhật không nắp có thể tích $ 200 $ m$ ^3 $. Đáy bể là hình chữ nhật có chiều dài gấp đôi chiều rộng. Giá thuê công nhân xây bể là $ 300~000 $ đồng/m$ ^2 $. Chi phí thuê công nhân thấp nhất là
% 	\choice
% 	{$ 36 $ triệu đồng}
% 	{\True $ 51 $ triệu đồng}
% 	{$ 75 $ triệu đồng}
% 	{$ 46 $ triệu đồng}
% 	\loigiai{
% 		Gọi chiều rộng, chiều dài của đáy lần lượt là $ x $ và $ 2x $, chiều cao là $ y $ $ (x, y>0) $.\\
% 		Diện tích các mặt bên và mặt đáy là $ S=6xy+2x^2 $.\\
% 		Thể tích là $ V=2x^2y=200\Rightarrow xy=\dfrac{100}{x}\cdot $\\ 
% 		Do đó, $ S=\dfrac{600}{x}+2x^2=\dfrac{300}{x}+\dfrac{300}{x}+2x^2\ge 3\sqrt[3]{\dfrac{300}{x}\cdot\dfrac{300}{x}\cdot 2x^2}=30\sqrt[3]{180} $.\\
% 		Vậy chi phí thấp nhất là $ T=30\sqrt[3]{180}\cdot 300~000\approx 51 $ triệu.}
% \end{ex}
% %%%=========HetCau_39=========%%%

% %%%=========Cau_40=========%%%
% \begin{ex}%[2H3K3-2]
% 	\renewcommand{\dotEX}{}
% 	Trong không gian $ Oxyz $, đường thẳng đi qua điểm $ M(1;2;2) $, song song với mặt phẳng $ (P)\colon x-y+z+3=0 $ đồng thời cắt đường thẳng $ d\colon\dfrac{x-1}{1}=\dfrac{y-2}{1}=\dfrac{z-3}{1} $ có phương trình là
% 	\choice
% 	{$ \left\{\begin{aligned}
% 	& x=1-t \\
% 	& y=2+t \\
% 	& z=2.
% \end{aligned}\right. $}
% 	{$ \left\{\begin{aligned}
% 	& x=1+t \\
% 	& y=2-t \\
% 	& z=2.
% \end{aligned}\right. $}
% 	{$ \left\{\begin{aligned}
% 	& x=1-t \\
% 	& y=2-t \\
% 	& z=2-t.
% \end{aligned}\right. $}
% 	{\True $ \left\{\begin{aligned}
% 	& x=1-t \\
% 	& y=2-t \\
% 	& z=2.
% \end{aligned}\right. $}
% 	\loigiai{
% 		Phương trình tham số của đường thẳng $ d\colon\left\{\begin{aligned}
% 	& x=1+t \\
% 	& y=2+t \\
% 	& z=3+t.
% \end{aligned}\right. $ \\
% 		Gọi $ \Delta $ là đường thẳng cần tìm. Theo đề bài, $ d $ cắt $ \Delta $ nên gọi $ I=\Delta \cap d\Rightarrow I\in d $ suy ra $ I(1+t;2+t;3+t) $.\\
% 		Ta có $ \overrightarrow{MI}=(t;t;t+1) $, mặt phẳng $ (P) $ có một véc-tơ pháp tuyến là $ \vec{n}=(1;-1;1) $.\\
% 		Đường thẳng $ \Delta $ song song với mặt phẳng $ (P) $ nên $ \begin{aligned}[t]
% 			\overrightarrow{MI}\perp \overrightarrow{n}&\Leftrightarrow \overrightarrow{MI}\cdot\overrightarrow{n}=0\\&\Leftrightarrow 1\cdot t+(-1)\cdot t+1\cdot(1+t)=0\\&\Leftrightarrow t=-1.
% 		\end{aligned} $ \\
% 		Suy ra $ \overrightarrow{MI}=(-1;-1;0) $ là một véc-tơ chỉ phương của đường thẳng $ \Delta $ và $ \Delta $ đi qua điểm $ M(1;2;2) $.\\
% 		Vậy phương trình tham số của đường thẳng $ \Delta $ là $ \left\{\begin{aligned}
% 	& x=1-t' \\
% 	& y=2-t' \\
% 	& z=2.
% \end{aligned}\right. $}
% \end{ex}
% %%%=========HetCau_40=========%%%

%%%=========Cau_41=========%%%
\begin{ex}%[2D4G5-2]
	\renewcommand{\dotEX}{.}
	Cho số phức $ z=a+bi $ $ (a, b\in \mathbb{R}) $ thỏa mãn $ | z |=1 $. Tìm giá trị lớn nhất của biểu thức $ A=| z+2 |+2| z-2 | $.
	\choice
	{$ 10\sqrt{2} $}
	{$ 7 $}
	{$ 10 $}
	{\True $ 5\sqrt{2} $}
	\loigiai{
		Ta có $ \left\{\begin{aligned}
	& |z+2|^2=(a+2)^2+b^2\\
	& |z-2|^2=(a-2)^2+b^2
\end{aligned}\right.\Rightarrow |z+2|^2+|z-2|^2=2(a^2+b^2)+8=2|z|^2+8=10 $.\\
		Mà $ A^2=(|z+2|+2|z-2|)^2\le (1^2+2^2)(|z+2|^2+|z-2|^2)=50 $.\\
		Vì $ A\ge 0 $ nên từ đó suy ra $ A\le \sqrt{50}=5\sqrt{2} $.\\
		Với $a=-\dfrac{3}{4}$; $b=\dfrac{\sqrt{7}}{4}$ ta có $A=5\sqrt{2}$.\\
		Vậy giá trị lớn nhất của $ A $ là $ 5\sqrt{2} $}
\end{ex}
%%%=========HetCau_41=========%%%

%%%=========Cau_42=========%%%
\begin{ex}%[2D3G2-4]
	Cho hàm số $ f(x) $ xác định và có đạo hàm $ f'(x) $ liên tục trên đoạn $ [1; 3] $ và $ f(x)\ne 0 $ với mọi $ x\in [1; 3] $, đồng thời $ f'(x)+[1+f(x)]^2=[f^2(x)(x-1)]^2 $ và $ f(1)=-1 $. Biết rằng $ \displaystyle\int\limits_1^3f(x)\mathrm{\,d}x=a\ln 3+b $ trong đó $ a, b\in \mathbb{Z} $. Tính tổng $ S=a+b^2 $.
	\choice
	{\True $ S=-1 $}
	{$ S=2 $}
	{$ S=0 $}
	{$ S=-4 $}
	\loigiai{
		Ta có $\begin{aligned}[t]
			f'(x)[1+f(x)]^2=[f^2(x)(x-1)]^2 \Leftrightarrow \dfrac{f'(x)[1+f(x)]^2}{f^4(x)}=(x-1)^2.
		\end{aligned}$ \\
		Lấy nguyên hàm hai vế ta được $$\begin{aligned}[t]
			\displaystyle\int{\dfrac{f'(x)[1+f(x)]^2}{f^4(x)}\mathrm{\,d}x=\displaystyle\int{(x-1)^2\mathrm{\,d}x}}&\Leftrightarrow \displaystyle\int{\dfrac{[1+2f(x)+f^2(x)]f'(x)}{f^4(x)}\mathrm{\,d}x}=\displaystyle\int{(x-1)^2\mathrm{\,d}x}\\&\Leftrightarrow \displaystyle\int\left[\dfrac{1}{f^4(x)}+\dfrac{2}{f^3(x)}+\dfrac{1}{f^2(x)}\right]\mathrm{\,d}[f(x)]=\dfrac{(x-1)^3}{3}+C\\&\Leftrightarrow -\dfrac{1}{3f^3(x)}-\dfrac{1}{f^2(x)}-\dfrac{1}{f(x)}=\dfrac{(x-1)^3}{3}+C\\&\Leftrightarrow -\dfrac{1+3f(x)+3f^2(x)}{3f^3(x)}=\dfrac{(x-1)^3}{3}+C
		\end{aligned}$$ 
		Mà $ f(1)=-1\Leftrightarrow-\dfrac{1-3+3}{-3}=C\Leftrightarrow C=\dfrac{1}{3}\cdot $\\
		Do đó $\begin{aligned}[t]
			-\dfrac{1+3f(x)+3f^2(x)}{3f^3(x)}=\dfrac{(x-1)^3}{3}+\dfrac{1}{3} &\Leftrightarrow \dfrac{1+3f(x)+3f^2(x)}{3f^3(x)}+\dfrac{1}{3}=-\dfrac{(x-1)^3}{3}\\
			&\Leftrightarrow \dfrac{[1+f(x)]^3}{f^3(x)}=-(x-1)^3\\&\Leftrightarrow \left[1+\dfrac{1}{f(x)}\right]^3=(1-x)^3\\ 
			&\Leftrightarrow f(x)=-\dfrac{1}{x}\cdot
		\end{aligned}$ \\
		Vậy $\displaystyle\int\limits_1^3f(x)\mathrm{\,d}x=\displaystyle\int\limits_1^3-\dfrac{1}{x}\mathrm{\,d}x=-\ln |x|\bigg|_1^3=-\ln 3$.\\ 
		Suy ra $a=-1$, $b=0$  $\Rightarrow S=a+b^2=-1$.
}
\end{ex}
%%%=========HetCau_42=========%%%

%%%=========Cau_43=========%%%
\begin{ex}%[2D2G6-5]
	Có bao nhiêu bộ $ (x; y) $? Với $ x $, $ y $ nguyên và $ 1\le x, y\le 2020 $ thỏa mãn \[(xy+2x+4y+8)\log_3\left(\dfrac{2y}{y+2}\right)\le (2x+3y-xy-6)\log_2\left(\dfrac{2x+1}{x-3}\right) . \] 
	\choice
	{\True $ 4\,034 $}
	{$ 2 $}
	{$ 2\,017 $}
	{$ 2\,017\cdot 2\,020 $}
	\loigiai{Điều kiện $ \left\{\begin{aligned}
	& x,y\in \mathbb{N}^\ast\colon x,y\le 2020 \\
	& \dfrac{2x+1}{x-3}>0,\, \dfrac{2y}{y+2}>0
\end{aligned}\right.\Leftrightarrow\left\{\begin{aligned}
	& x, y\in \mathbb{N}^\ast\colon x,y\le 2020 \\
	& x>3,\, y>0.
\end{aligned}\right. $ \\ 
		Bất phương trình đã cho được viết lại \[(x-3)(y-2)\log_2\left(\dfrac{x+4}{x-3}+1\right)+(x+4)(y+2)\log_3\left(\dfrac{y-2}{y+2}+1\right)\le 0. \,(*) \] 
		\begin{itemize}
			\item Với $ y=1 $ thì $ (*) $ trở thành $ -(x-3)\log_2\left(\dfrac{x+4}{x-3}+1\right)+3(x+4)\log_3\dfrac{2}{3}\le 0 $, rõ ràng bất phương trình này nghiệm đúng với mọi $ x>3 $ vì $ -(x-3)<0 $, $ \log_2\left(\dfrac{x+4}{x-3}+1\right)>\log_2(0+1)=0 $, $ 3(x+4)>0 $, $ \log_3\dfrac{2}{3}<0 $.\\
			Như vậy trường hợp này cho ta đúng $ 2\,017 $ bộ $ (x;y)=(x;1) $ với $ 4\le x\le 2020,x\in \mathbb{N} $.
			\item Với $ y=2 $ thì $ (*) $ thành $ 4(x+4)\log_31\le 0 $, bất phương trình này cũng luôn đúng với mọi $ x $ mà $ 4\le x\le 2020 $, $ x\in \mathbb{N} $. Trường hợp này cho ta thêm $ 2\,017 $ cặp $ (x;y) $.
			\item Với $ y>2 $, $ x>3 $ thì vế trái của bất phương trình $ (*) $ luôn dương nên $ (*) $ không xảy ra.
		\end{itemize}
		Vậy có đúng $ 4\,034 $ bộ số $ (x;y) $ thỏa mãn yêu cầu bài toán.}
\end{ex}
%%%=========HetCau_43=========%%%

%%%=========Cau_44=========%%%
\begin{ex}%[1H3K4-3]
	\renewcommand{\dotEX}{}
	\immini{Cho hình lăng trụ đều $ ABC. A'B'C' $ có tất cả các cạnh bằng nhau và bằng $ 2a $ (minh họa như hình vẽ). Cô-sin của góc hợp bởi $ (A'BC) $ và $ (ABC) $ bằng
		\choice
		{$ \dfrac{\sqrt{21}}{3}\cdot $}
		{\True $ \dfrac{\sqrt{21}}{7}\cdot $}
		{$ \dfrac{2}{\sqrt{3}}\cdot $}
		{$ \dfrac{2}{\sqrt{7}}\cdot $}}{
		\begin{tikzpicture}[>=stealth,line join=round,line cap=round,font=\scriptsize,scale=0.65]
			\coordinate [label=below left: $ A $] (A) at (0,0);
			\coordinate [label=below right: $ B $] (B) at (4,0);
			\coordinate [label=below: $ C $] (C) at (3,-1);
			\coordinate [label=above left: $ A' $](A') at ($ (A)+(0,3) $);
			\coordinate [label=below right: $ C' $](C') at ($ (A')+(C)-(A) $);
			\coordinate [label=above right: $ B' $](B') at ($ (A')+(B)-(A) $);
			\draw (C')--(A')--(A)--(C)--(B)--(B')--(C')--(C)--(A') (A')--(B');
			\draw [dashed] (A)--(B)--(A');
			\draw [fill=black] (A)circle(1pt) (B)circle(1pt) (C)circle(1pt) (A')circle(1pt) (B')circle(1pt) (C')circle(1pt);
		\end{tikzpicture}
	}
	\loigiai{
		\immini{Gọi $ I $ là trung điểm của $ BC $.\\ Khi đó, $ \left\{\begin{aligned}
	& BC\perp AI\\
	& BC\perp AA'
\end{aligned}\right. $ nên $ BC\perp (AA'I)\Rightarrow BC\perp A'I $.\\ Vậy góc hợp bởi $ (A'BC) $ và $ (ABC) $ bằng $ \widehat{AIA'} $.\\
			Ta có $ AI=\dfrac{2a\sqrt{3}}{2}=a\sqrt{3} $, $ AA'=2a\Rightarrow \tan \widehat{AIA'}=\dfrac{AA'}{AI}=\dfrac{2a}{a\sqrt{3}}=\dfrac{2}{\sqrt{3}}\cdot $
		}{
			\begin{tikzpicture}[>=stealth,line join=round,line cap=round,font=\scriptsize,scale=0.65]
				\coordinate [label=below left: $ A $] (A) at (0,0);
				\coordinate [label=below right: $ B $] (B) at (4,0);
				\coordinate [label=below: $ C $] (C) at (3,-1);
				\coordinate [label=above left: $ A' $](A') at ($ (A)+(0,3) $);
				\coordinate [label=below right: $ C' $](C') at ($ (A')+(C)-(A) $);
				\coordinate [label=above right: $ B' $](B') at ($ (A')+(B)-(A) $);
				\coordinate [label=below right: $ I $](I) at ($ (C)!0.5!(B) $);
				\draw (C')--(A')--(A)--(C)--(B)--(B')--(C')--(C)--(A') (A')--(B');
				\draw [dashed] (A)--(B)--(A')--(I)--(A);
				\draw [fill=black] (A)circle(1pt) (B)circle(1pt) (C)circle(1pt) (A')circle(1pt) (B')circle(1pt) (C')circle(1pt) (I)circle(1pt);				
			\end{tikzpicture}
		}
		\noindent	Mặt khác $ \begin{aligned}[t]
			1+\tan^2\widehat{AIA'}=\dfrac{1}{\cos^2\widehat{AIA'}}\Rightarrow {\cos^2}\widehat{AIA'}=\dfrac{1}{1+\tan^2\widehat{AIA'}}=\dfrac{1}{1+\dfrac{4}{3}}=\dfrac{3}{7}\cdot
		\end{aligned} $ \\
		Vậy $ \cos \widehat{AIA'}=\dfrac{\sqrt{21}}{7}\cdot $
	}
\end{ex}
%%%=========HetCau_44=========%%%

%%%=========Cau_45=========%%%
\begin{ex}%[2H1K3-2]
	\renewcommand{\dotEX}{}
	Cho hình chóp $ S. ABC $ có đáy $ ABC $ là tam giác đều, $ SA\perp (ABC) $. Mặt phẳng $ (SBC) $ cách $ A $ một khoảng bằng $ a $ và hợp với mặt phẳng $ (ABC) $ góc $ 30^\circ $. Thể tích của khối chóp $ S. ABC $ bằng	
	\choice
	{\True $ \dfrac{8a^3}{9}\cdot $}
	{$ \dfrac{\sqrt{3}a^3}{12}\cdot $}
	{$ \dfrac{4a^3}{9}\cdot $}
	{$ \dfrac{8a^3}{3}\cdot $}
	\loigiai{
		\immini{Gọi $ I $ là trung điểm $ BC $ suy ra góc giữa $ (SBC) $ và $ (ABC) $ là $ \widehat{SIA}=30^\circ $.\\ 
			Gọi $ H $ là hình chiếu của $ A $ trên $ SI \Rightarrow AH\perp (SBC)$ suy ra $ \mathrm{d}(A,(SBC))=AH=a $.\\
			Xét tam giác $ AHI $ vuông tại $ H $, ta có $ AI=\dfrac{AH}{\sin{30^\circ}}=2a $.\\
			Giả sử tam giác đều $ ABC $ có cạnh bằng $ x $, mà $ AI $ là đường cao suy ra $ 2a=x\cdot\dfrac{\sqrt{3}}{2}\Rightarrow x=\dfrac{4a}{\sqrt{3}}\cdot $\\
			Diện tích tam giác đều $ ABC $ là $ S_{ABC}=\left(\dfrac{4a}{\sqrt{3}}\right)^2\cdot\dfrac{\sqrt{3}}{4}=\dfrac{4a^2\sqrt{3}}{3}\cdot $\\
			Xét tam giác $ SAI $ vuông tại $ A $ suy ra $ SA=AI\cdot\tan{30^\circ}=\dfrac{2a}{\sqrt{3}}\cdot $\\
			Vậy $ V_{S. ABC}=\dfrac{1}{3}\cdot S_{ABC}\cdot SA=\dfrac{1}{3}\cdot\dfrac{4a^2\sqrt{3}}{3}\cdot\dfrac{2a}{\sqrt{3}}=\dfrac{8a^3}{9}\cdot $}{
			\begin{tikzpicture}[>=stealth,line join=round,line cap=round,font=\scriptsize,scale=0.75]
				\coordinate [label=left: $ A $] (A) at (0,0);
				\coordinate [label=below: $ B $] (B) at (1.5,-1);
				\coordinate [label=right: $ C $] (C) at (4,0);
				\coordinate [label=above: $ S $](S) at ($ (A)+(0,3) $);
				\coordinate [label=below: $ I $](I) at ($ (C)!0.5!(B) $);
				\coordinate [label=right: $ H $](H) at ($ (S)!0.4!(I) $);
				\draw (I)--(S)--(A)--(B)--(C)--(S)--(B);
				\draw [dashed] (I)--(A)--(H) (A)--(C);
				\draw [fill=black] (A)circle(1pt) (B)circle(1pt) (C)circle(1pt) (S)circle(1pt);
				\foreach \x/\dinh/\y in {A/I/B,A/H/S}\draw[fill = gray!50] ($ (\dinh)!5pt!(\x) $)--($ (\dinh)!5pt!(\x)+(\dinh)!5pt!(\y)-(\dinh) $)--($ (\dinh)!5pt!(\y) $)--(\dinh)--cycle; 
				\draw [fill=black]  (H)circle(1pt)   (I)circle(1pt);
			\end{tikzpicture}
		}
	}
\end{ex}
%%%=========HetCau_45=========%%%

%%%=========Cau_46=========%%%
\begin{ex}%[2D1G3-6]
	\immini{Cho hàm số $ f(x) $ liên tục trên $ \mathbb{R}, $ có đồ thị như hình vẽ.
		Có tất cả bao nhiêu giá trị nguyên của tham số $ a $ để hàm số $ y=\left| f\left(\dfrac{8x}{x^2+1}\right)+a-1 \right| $ có giá trị lớn nhất không vượt quá $ 20 $?
		\choice
		{$ 41 $}
		{\True $ 31 $}
		{$ 35 $}
		{$ 29 $}}{
		\begin{tikzpicture}[>=stealth,line join=round,line cap=round,font=\scriptsize,scale=0.5]
			\draw [->](-6,0)--(5,0)node[below]{$ x $};
			\draw [->] (0,-5)--(0,7.5)node[left]{$ y $};
			\draw [smooth,domain=-6:4.07] plot(\x,{15/472*(\x)^4+10/59*(\x)^3-15/59*(\x)^2-120/59*(\x)-46/59});
			\draw [dashed] (-4,0)node[below]{$ -4 $}--(-4,34/59)--(0,34/59) (-2,0)node[below]{$ -2 $}--(-2,84/59)--(0,84/59) (2,0)node[above]{$ 2 $}--(2,-4)--(0,-4)node[left]{$ -4 $} (4,0)node[below]{$ 4 $}--(4,6)--(0,6)node[left]{$ 6 $};
			\draw [fill=black] (-4,34/59)circle(1pt) (-2,84/59)circle(1pt) (2,-4)circle(1pt) (4,6)circle(1pt) (0,2)node[left]{$ 2 $}circle(1pt) (0,4)node[left]{$ 4 $}circle(1pt) (0,0)node[below right]{$ O $}circle(1pt);
		\end{tikzpicture}
	}
	\loigiai{
		Đặt $ t=\dfrac{8x}{x^2+1}\cdot $\\
		Ta có $ t'=\dfrac{-8x^2+8}{(x^2+1)^2};\,t'=0\Leftrightarrow x=\pm 1 $.\\
		Bảng biến thiên của $ t(x) $ 
		\begin{center}
			\begin{tikzpicture}
				\tkzTabInit[nocadre,lgt=1.2,espcl=2.5,deltacl=0.6]
				{$ x $ /0.6, $ t' $ /0.6, $ t $ /2}
				{$ -\infty $, $ -1 $, $ 1 $, $ +\infty $}
				\tkzTabLine{,-,0,+,0,-,}
				\tkzTabVar{+/ $ 0 $,-/ $ -4 $,+/ $ 4 $,-/ $ 0 $}
			\end{tikzpicture}
		\end{center}
		Dựa vào bảng biến thiên, suy ra $ t\in [-4;4 ] $.\\
		Xét hàm số $ h(t)=f(t)+a-1,t\in [-4;4 ], $ ta có \[h'(t)=f'(t), h'(t)=0\Leftrightarrow f'(t)=0\Leftrightarrow \left[\begin{aligned}
	& t=-4\in [-4;4 ] \\
	& t=-2\in [-4;4 ] \\
	& t=2\in [-4;4 ].
\end{aligned}\right. \]
		Suy ra $ \max\limits_{[-4;4 ]}| h(t) |=\max\{| a+5 |;| a-5 | \} $.\\
		Yêu cầu bài toán $ \Leftrightarrow \left\{\begin{aligned}
	& | a+5 |\le 20 \\
	& | a-5 |\le 20
\end{aligned}\right.\Leftrightarrow \left\{\begin{aligned}
	& -20\le a+5\le 20 \\
	& -20\le a-5\le 20
\end{aligned}\right.\Leftrightarrow \left\{\begin{aligned}
	& -25\le a\le 15 \\
	& -15\le a\le 25
\end{aligned}\right.\Leftrightarrow -15\le a\le 15 $.\\
		Vậy có tất cả $ 31 $ giá trị nguyên của tham số $ a $ thỏa mãn yêu cầu bài toán.
	}
\end{ex}
%%%=========HetCau_46=========%%%

%%%=========Cau_47=========%%%
\begin{ex}%[2D3G3-1]
	\renewcommand{\dotEX}{}
	\immini{Cho $ f(x) $ là hàm đa thức bậc $ 3 $ có đồ thị như hình vẽ. Tiếp tuyến của đồ thị hàm số tại điểm $ M $ có hoành độ bằng $ -2 $ cắt đồ thị tại điểm thứ hai $ N(1;1) $ và cắt $ Ox $ tại điểm có hoành độ bằng $ 4 $. Biết diện tích phần gạch chéo là $ \dfrac{9}{16} $. Tích phân $ \displaystyle\int\limits_{-1}^1f(x)\mathrm{\,d}x $ bằng
		\choice
		{$ \dfrac{31}{18}\cdot $}
		{\True $ \dfrac{13}{6}\cdot $}
		{$ \dfrac{19}{9}\cdot $}
		{$ \dfrac{7}{3}\cdot $}}{
		\begin{tikzpicture}[>=stealth,line join=round,line cap=round,font=\scriptsize,scale=0.65]
			\draw [->] (-4,0)--(5,0)node[below]{$ x $};
			\draw [->] (0,-1)--(0,4)node[left]{$ y $};
			\draw [smooth,domain=-4:2.7] plot(\x,{(1/12)*(\x)^3+(1/4)*(\x)^2-1/3*(\x)+1});
			\draw [smooth,domain=-4:4.6] plot(\x,{-1/3*(\x)+4/3});
			\draw [fill=black] (-2,2)circle(1pt) (1,1)circle(1pt) (4,0)circle(1pt) (0,0)node[below left]{$ O $}circle(1pt);
			\draw [dashed] (-2,0)node[below]{$ -2 $}--(-2,2)--(0,2)node[right]{$ 2 $} (1,0)node[below]{$ 1 $}--(1,1)--(0,1)node[below left]{$ 1 $} (4,0)node[below]{$ 4 $};
			\draw [pattern = north west lines, domain=-2:1, samples=100] plot(\x,{(1/12)*(\x)^3+(1/4)*(\x)^2-1/3*(\x)+1})--(1,1)--(-2,2);
		\end{tikzpicture}
	}
	\loigiai{
		Dựa vào giả thiết đường thẳng $ d $ đi qua hai điểm $ M(-2;2) $ và $ P(4;0) $.\\ Suy ra $ d\colon x+3y-4=0\Rightarrow y=-\dfrac{1}{3}x+\dfrac{4}{3}\cdot $\\ 
		Từ giả thiết, ta có hàm số $ f(x)=ax^3+bx^2+cx+1\Rightarrow f'(x)=3ax^2+2bx+c $.\\
		 Đồ thị hàm số lần lượt qua điểm $(-2;2)$; qua điểm $(1;1)$ và tiếp xúc đường thẳng $ d $ tại $ x=-2 $ 
		 $ \begin{aligned}[t]
			&\Leftrightarrow\left\{\begin{aligned}
	& 1=-8a+4b-2c \\
	& 0=a+b+c \\
	& 12a-4b+c=-\dfrac{1}{3}
\end{aligned}\right.\\&\Rightarrow \left\{\begin{aligned}
	& a=\dfrac{1}{12}\\
	& b=\dfrac{1}{4}\\
	& c=-\dfrac{1}{3}
\end{aligned}\right.\Rightarrow y=\dfrac{1}{12}x^3+\dfrac{1}{4}x^2-\dfrac{1}{3}x+1.
		\end{aligned} $ \\
		Từ đó $ \displaystyle\int\limits_{-1}^1f(x)\mathrm{\,d}x=\dfrac{13}{6}\cdot $}
\end{ex}
%%%=========HetCau_47=========%%%

%%%=========Cau_48=========%%%
\begin{ex}%[2D2G5-5]
	Tổng tất cả các giá trị của tham số $ m $ để phương trình \[3^{x^2-2x+1-2| x-m |}=\log_{x^2-2x+3}(2| x-m |+2) \] có đúng ba nghiệm phân biệt là
	\choice
	{\True $ 3 $}
	{$ 0 $}
	{$ 2 $}
	{$ 1 $}
	\loigiai{
		Ta có $ \begin{aligned}[t]
			&3^{x^2-2x+1-2| x-m |}=\log_{x^2-2x+3}(2| x-m |+2)\\ \Leftrightarrow\, &3^{x^2-2x+3-(2| x-m |+2)}=\dfrac{\ln (2| x-m |+2)}{\ln (x^2-2x+3)}\\\ 
 \Leftrightarrow\, &{3^{x^2-2x+3}}\cdot\ln (x^2-2x+3)=3^{2| x-m |+2}\cdot\ln (2| x-m |+2). \, (*)
		\end{aligned} $ \\
		Xét hàm đặc trưng $ f(t)=3^t\cdot\ln t $, $ t\ge 2 $ là hàm số đồng biến.\\ Suy ra
		 $ (*)\Leftrightarrow x^2-2x+3=2| x-m |+2\Leftrightarrow x^2-2x-2| x-m |+1=0 $.\\ Đặt $ g(x)= x^2-2x-2| x-m |+1 $, ta có
		$g(x)=\left\{\begin{aligned}
	& x^2-4x+2m+2\ \ \text{khi}\;x\ge m \\
	& x^2-2m+1\qquad\quad\text{khi}\; x\le m
\end{aligned}\right.;$
\[ g'(x)=\left\{\begin{aligned}
	& 2x-4\ \ \text{khi}\; x> m \\
	& 2x\qquad\text{khi}\; x< m
\end{aligned}\right.; \ g'(x)=0\Leftrightarrow \left[\begin{aligned}
	& x=2\;\; \text{khi}\; x> m \\
	& x=0\;\;\text{khi}\; x< m.
\end{aligned}\right.\]
		\begin{itemize}
			\item \textbf{Trường hợp 1.} $ m\le 0 $ ta có bảng biến thiên của $ g(x) $ như sau
			\begin{center}
				\begin{tikzpicture}
					\tkzTabInit[nocadre,lgt=1.2,espcl=2.5,deltacl=0.6]
					{$ x $ /0.6, $ g'(x) $ /0.6, $ g(x) $ /2}
					{$ -\infty $, $ m $, $ 0 $, $ 2 $, $ +\infty $}
					\tkzTabLine{,-,,-,0,-,0,+,}
					\tkzTabVar{+/ $ +\infty $,R,R,-/ $ g(2) $,+/ $ +\infty $}
				\end{tikzpicture}
			\end{center}
			Phương trình chỉ có tối đa $ 2 $ nghiệm nên không có $ m $ thỏa mãn.
			\item \textbf{Trường hợp 2.} $ m\ge 2 $ tương tự không có $m$ thỏa mãn.
			\item \textbf{Trường hợp 3.} $ 0<m<2, $ bảng biến thiên $ g(x) $ như sau
			\begin{center}
				\begin{tikzpicture}
					\tkzTabInit[nocadre,lgt=1.2,espcl=2.5,deltacl=0.6]
					{$ x $ /0.6, $ g'(x) $ /0.6, $ g(x) $ /2}
					{$ -\infty $, $ 0 $, $ m $, $ 2 $, $ +\infty $}
					\tkzTabLine{,-,0,+,,-,0,+,}
					\tkzTabVar{+/ $ +\infty $,-/ $ -2m+1 $,+/ $ (m-1)^2 $,-/ $ 2m-3 $,+/ $ +\infty $}
				\end{tikzpicture}
			\end{center}
			Phương trình có $ 3 $ nghiệm khi $ \left[\begin{aligned}
	& (m-1)^2=0 \\
	& -2m+1=0>2m-3 \\
	& -2m+1<0=2m-3
\end{aligned}\right.\Leftrightarrow \left[\begin{aligned}
	& m=1 \\
	& m=\dfrac{1}{2}\\
	& m=\dfrac{3}{2}\cdot
\end{aligned}\right. $ 
		\end{itemize}
	}
\end{ex}
%%%=========HetCau_48=========%%%

%%%=========Cau_49=========%%%
\begin{ex}%[2D4G1-2]
	\renewcommand{\dotEX}{}
	Cho các số phức $ z_1=1+3i,z_2=-5-3i $. Tìm điểm $ M(x;y) $ biểu diễn số phức $ z_3 $, biết rằng trong mặt phẳng phức điểm $ M $ nằm trên đường thẳng $ x-2y+1=0 $ và mô-đun số phức $ \mathrm{w}=3z_3-z_2-2z_1 $ đạt giá trị nhỏ nhất.
	\choice
	{$ M\left(\dfrac{3}{5};\dfrac{1}{5}\right)\cdot $}
	{$ M\left(-\dfrac{3}{5};-\dfrac{1}{5}\right)\cdot $}
	{$ M\left(\dfrac{3}{5};-\dfrac{1}{5}\right)\cdot $}
	{\True $ M\left(-\dfrac{3}{5};\dfrac{1}{5}\right)\cdot $}
	\loigiai{
		\textbf{Trắc nghiệm.} Thay tọa độ điểm $ M\left(-\dfrac{3}{5};\dfrac{1}{5}\right) $ vào vế trái phương trình đường thẳng kết quả bằng $ 0 $ thỏa ta chọn $ M\left(-\dfrac{3}{5};\dfrac{1}{5}\right)$.\\
		\textbf{Tự luận.}
		Ta có $ w=3z_3-z_2-2z_1=3z_3+3-3i=3(z_3+1-i) $, suy ra $ | w |=3| z_3+1-i |=3AM $ với $ A(-1;1) $.\\
		Điểm $ M(x;y) $ biểu diễn số phức $ z_3 $ nằm trên đường thẳng $ d\colon x-2y+1=0 $ và $ A(-1;1)\notin d $.\\ 
		Khi đó $ | w |=3| z_3+1-i |=3AM $ đạt giá trị nhỏ nhất khi $ AM $ ngắn nhất $ \Leftrightarrow AM\perp d $.\\
		Vì $ AM\perp d $ nên $ AM\colon 2x+y+1=0 $. 
		Khi đó $ M=AM\cap d $ nên $ M\left(-\dfrac{3}{5};\dfrac{1}{5}\right)$.}
\end{ex}
%%%=========HetCau_49=========%%%

%%%=========Cau_50=========%%%
\begin{ex}%[2H3G2-8]
	Trong không gian $ Oxyz $, cho ba điểm $ A(2;-2;4) $, $ B(-3;3;-1) $, $ C(-1;-1;-1) $ và mặt phẳng $ (P)\colon 2x-y+2z+8=0 $. Xét điểm $ M $ thay đổi thuộc $ (P) $, tìm giá trị nhỏ nhất của biểu thức $ T=2MA^2+MB^2-MC^2 $.
	\choice
	{\True $ 102 $}
	{$ 35 $}
	{$ 105 $}
	{$ 30 $}
	\loigiai{
		Gọi $ I $ là điểm thỏa mãn  \begin{align*}
			2\overrightarrow{IA}+\overrightarrow{IB}-\overrightarrow{IC}=\vec{0}&\Leftrightarrow 2 (\overrightarrow{OA}-\overrightarrow{OI})+ (\overrightarrow{OB}-\overrightarrow{OI})- (\overrightarrow{OC}-\overrightarrow{OI})=\overrightarrow{0}\\&\Leftrightarrow \overrightarrow{OI}=\overrightarrow{OA}+\dfrac{1}{2}\overrightarrow{OB}-\dfrac{1}{2}\overrightarrow{OC}=(1;0;4)\\&\Leftrightarrow I(1;0;4).
		\end{align*}  \\
		Khi đó, với mọi điểm $ M(x;y;z)\in (P) $, ta luôn có $$\begin{aligned}[t]
			T&=2 (\overrightarrow{MI}+\overrightarrow{IA})^2+ (\overrightarrow{MI}+\overrightarrow{IB})^2- (\overrightarrow{MI}+\overrightarrow{IC})^2\\&=2\overrightarrow{MI}^2+2\overrightarrow{MI}\cdot(2\overrightarrow{IA}+\overrightarrow{IB}-\overrightarrow{IC})+2\overrightarrow{IA}^2+\overrightarrow{IB}^2-\overrightarrow{IC}^2\\&=2MI^2+2IA^2+IB^2-IC^2.
		\end{aligned}$$ 
		Ta tính được $ 2IA^2+IB^2-IC^2=30 $.\\
		Do đó, $ T $ đạt giá trị nhỏ nhất $ \Leftrightarrow MI $ đạt giá trị nhỏ nhất $ \Leftrightarrow MI\perp (P) $.\\
		Lúc này, $ IM=\mathrm{d}(I,(P))=\dfrac{| 2\cdot 1-0+2\cdot4+8 |}{\sqrt{2^2+(-1)^2+2^2}}=6 $.\\
		Vậy $ T_{\min}=2\cdot6^2+30=102 $.}
\end{ex}
%%%=========HetCau_50=========%%%


\Closesolutionfile{ans}